\textbf{Время:} 01:45

\textbf{Тезис:} время дрейфа электронов восстанавливает координату трека, а не скорость пролетающей частицы.

\textbf{Визуал:} частица → ионизационные электроны → дрейф к торцу → считывающие площадки.

\textbf{Текст:} Чтобы измерить обещанную C, сначала нужны надёжно восстановленные импульсы. Начнём с трекинга в газовой камере. Заряженная частица ионизует атомы: вдоль её пути остаются электроны. По ним нужно восстановить пространственный след.

Электрическое поле переносит электроны к считывающим плоскостям на торцах камеры. В двух половинах TPC дрейф направлен к противоположным торцам. Положение сработавших площадок даёт две координаты, а время прихода заряда — третью. Если скорость дрейфа известна, произведение этой скорости на время дрейфа даёт расстояние от места ионизации до плоскости считывания.

Здесь важно различить две скорости. Пион уже пролетел через газ; мы ждём созданные им электроны. Время их дрейфа говорит о координате следа, а не о скорости пиона. Собрав множество таких пространственных точек и связав их в последовательность, восстанавливаем трек.

Пауза: пройти указкой один участок — от места ионизации до сигнала на торце.

\textbf{Переход:} По треку теперь можно определить импульс. Как именно это делается?
